% =====================================================================
% IEEE Transactions on Quantum Engineering (TQE) manuscript
% Built on the official TQE LaTeX template (ieeeaccess.cls, March 2020).
%
% Project layout
%   main.tex          this file: front matter and section skeleton
%   bibliography.tex  reference list (thebibliography, IEEE numbered style)
%   biography.tex     author biographies (IEEEbiography blocks)
%   supplementary.tex standalone Supplementary Information document
%   ../figures/       vector PDF figures shared with the original draft
%   photos/           author photographs (1 in x 1.25 in, 300 dpi)
%
% Compile: pdflatex main.tex (twice; no BibTeX step is required)
% =====================================================================
\documentclass{ieeeaccess}

\usepackage{cite}
\usepackage{amsmath,amssymb,amsfonts}
\usepackage{algorithmic}
\usepackage{graphicx}
\usepackage{textcomp}
\usepackage{booktabs}
\usepackage{multirow}
\usepackage{url}

\graphicspath{{../figures/}{./photos/}{./}}

% ---------------- Theorem environments (class-native \newtheorem; amsthm conflicts with ieeeaccess) --------
\newtheorem{theorem}{Theorem}
\newtheorem{proposition}{Proposition}
\newtheorem{corollary}{Corollary}
\newtheorem{lemma}{Lemma}
\newtheorem{definition}{Definition}
\newtheorem{remark}{Remark}
\newenvironment{proof}{\noindent\textit{Proof:}\ }{\hfill$\blacksquare$\par}

% ---------------- Notation macros ----------------
\newcommand{\B}{\mathbb{B}}
\newcommand{\R}{\mathbb{R}}
\newcommand{\PR}{P_{\mathrm{R}}}
\newcommand{\Hcost}{H_{F}}
\newcommand{\Hmix}{H_{\mathrm{M}}}
\newcommand{\Jres}{J_{\mathrm{res}}}
\newcommand{\Efib}{E_{\mathrm{fib}}}
\newcommand{\norminf}[1]{\lVert #1 \rVert_{\infty}}
\DeclareMathOperator*{\argmin}{arg\,min}
\DeclareMathOperator{\msupp}{msupp}

\def\BibTeX{{\rm B\kern-.05em{\sc i\kern-.025em b}\kern-.08em
    T\kern-.1667em\lower.7ex\hbox{E}\kern-.125emX}}

\begin{document}

\history{Date of publication xxxx 00, 0000, date of current version xxxx 00, 0000.}
\doi{10.1109/TQE.2026.DOI}

% ---------------- Title (working title; to be finalised) ----------------
\title{Exact Selective Locality Reduction for Resource-Constrained Quantum Approximate Optimization}

% ---------------- Authors and affiliations ----------------
\author{\uppercase{Jingcong Liu}\authorrefmark{1},
\uppercase{Yiding Su}\authorrefmark{1},
\uppercase{Ruojun Zhang}\authorrefmark{2},
\uppercase{Hafiz Muhammad Waseem}\authorrefmark{3},
\uppercase{Bo Chen}\authorrefmark{4},
and \uppercase{Carsten Maple}\authorrefmark{3}}

\address[1]{Department of Mathematics, University of Warwick, Coventry CV4 7AL,
United Kingdom}
\address[2]{School of Electronics and Computer Science, University of Southampton,
Southampton SO17 1BJ, United Kingdom}
\address[3]{Warwick Manufacturing Group, University of Warwick, Coventry CV4 7AL,
United Kingdom}
\address[4]{Warwick Business School, University of Warwick, Coventry CV4 7AL,
United Kingdom}

\tfootnote{This work was supported in part by the National Hub for Edge AI,
UK Research and Innovation (UKRI), through Research Grant EP/Y028813/1, in part by
the UKRI Cross Council under Grant UKRI3734, in part by the Academic Centre of
Excellence in Cyber Security Research, University of Warwick, under Grant
EP/R007195/1, and in part by the University of Warwick Undergraduate Research
Support Scheme (URSS).}

\markboth
{Liu \headeretal: Exact Selective Locality Reduction for Resource-Constrained QAOA}
{Liu \headeretal: Exact Selective Locality Reduction for Resource-Constrained QAOA}

\corresp{Corresponding author: Hafiz Muhammad Waseem (e-mail: HM.Waseem@warwick.ac.uk).}

% ---------------- Abstract (150 to 250 words, one paragraph, no citations) --------------
\begin{abstract}
Abstract to be drafted. The abstract must be a single paragraph of 150 to 250 words,
self-contained, without abbreviations, footnotes, references, displayed equations, or
tables, and should include three or four keywords or phrases.
\end{abstract}

\begin{keywords}
Pseudo-Boolean optimization, quadratization, quantum approximate optimization algorithm,
quantum circuit compilation, resource-aware quantum algorithms, warm-start QAOA.
\end{keywords}

\titlepgskip=-15pt

\maketitle

% =====================================================================
\section{Introduction}
\label{sec:introduction}
\PARstart{H}{igher-order} binary interactions arise naturally in clause penalties,
hypergraph objectives, and many-body spin models. Text to follow.

% Paragraph plan:
% (1) Higher-order objectives in QAOA and why representation is a design choice.
% (2) The two endpoint strategies and what each costs.
% (3) The gap: the exact intermediate family as a design space under joint budgets.
% (4) Contributions (numbered) and the main negative finding, stated plainly.

% =====================================================================
\section{Related Work}
\label{sec:related}

\subsection{Quadratization of Pseudo-Boolean Functions}
\label{sec:related:quadratization}

\subsection{Higher-Order Encodings and QAOA Circuit Synthesis}
\label{sec:related:higherorder}

\subsection{Hardware-Aware Reformulation and Compilation}
\label{sec:related:hardware}

\subsection{Warm Starts and Initialization}
\label{sec:related:warmstart}

% Table I: closest-work comparison (native, full exact quadratization,
% hardware-aware auxiliary selection, approximate quadratization, this work).
\begin{table}[!t]
\caption{Comparison of the proposed framework with the closest prior representation
strategies (to be completed).}
\label{tab:closest_work}
\centering
\scriptsize
\setlength{\tabcolsep}{3pt}
\begin{tabular}{lccccc}
\toprule
Strategy & Exact & Aux. & Locality & Connectivity & Selection \\
\midrule
Native higher-order & \checkmark & 0 & $\deg f$ & -- & -- \\
Full exact quadratization & \checkmark & max & 2 & -- & count \\
Hardware-aware aux.\ selection & \checkmark & max & 2 & \checkmark & degree \\
Approximate quadratization & -- & var. & 2 & -- & fidelity \\
This work & \checkmark & var. & mixed & \checkmark & multi-obj. \\
\bottomrule
\end{tabular}
\end{table}

% =====================================================================
% =====================================================================
\section{Exact Selective Locality Reduction}
\label{sec:reduction}

This section develops the exact representation family that lies between native
higher-order encoding and complete quadratization. We first fix notation, then state
a sharp penalty condition for one shared substitution, extend it to sequences of
substitutions, and specialise the result to cubic objectives. Table~\ref{tab:notation}
collects the symbols used throughout the paper.

Let $\B=\{0,1\}$, $[n]=\{1,\dots,n\}$, and $x_A=\prod_{i\in A}x_i$ with $x_\emptyset=1$.
Every function $f\colon\B^n\to\R$ has a unique multilinear expansion
$f(x)=\sum_{A\subseteq[n]}c_A x_A$, and its monomial support is
$\msupp(f)=\{A : c_A\neq 0\}$ with degree $\deg f=\max_{A\in\msupp(f)}|A|$
\cite{BorosHammer2002,Hadfield2021}. The diagonal encoding
$H_f=f(N_1,\dots,N_n)$ with $N_i=(I-Z_i)/2$ satisfies $H_f\ket{x}=f(x)\ket{x}$, and
its largest Pauli weight equals $\deg f$ \cite{Hadfield2021}. A representation
$F_\sigma\colon\B^{n}\times\B^{a_\sigma}\to\R$ with $a_\sigma$ auxiliary bits is
\emph{pointwise exact} when
\begin{equation}
f(x)=\min_{y\in\B^{a_\sigma}}F_\sigma(x,y)\quad\text{for every }x\in\B^n .
\label{eq:pointwise}
\end{equation}
This identity preserves the optimal value and the projected minimisers, but it leaves
the spectrum of $H_{F_\sigma}$ outside the optimal auxiliary fibres unconstrained
\cite{BorosGruber2014,Dattani2019}. The classical exactness of a representation
therefore does not determine its finite-depth QAOA behaviour, a distinction that
organises the experiments in Section~\ref{sec:experiments}. All constructions use the
Rosenberg penalty \cite{Rosenberg1975}
\begin{equation}
\PR(a,b,y)=ab-2ay-2by+3y,
\label{eq:rosenberg}
\end{equation}
which vanishes when $y=ab$ and is at least one otherwise.

\begin{table}[!t]
\caption{Core notation. Subscript $\sigma$ marks dependence on the chosen representation.}
\label{tab:notation}
\centering
\footnotesize
\setlength{\tabcolsep}{4pt}
\begin{tabular}{ll}
\toprule
Symbol & Meaning \\
\midrule
$x$, $y$, $u=(x,y)$ & Original bits, auxiliary bits, combined register \\
$x_A$, $c_A$ & Monomial $\prod_{i\in A}x_i$ and its coefficient in $f$ \\
$f$, $F_\sigma$ & Original objective and its pointwise-exact representation \\
$\sigma$, $a_\sigma$, $N_\sigma$ & Representation, auxiliary count, total qubits $n+a_\sigma$ \\
$D$, $R$ & Higher-order terms reduced and terms retained natively \\
$p$, $y_p$, $P_\rho$ & Variable pair, its product auxiliary, active pair set \\
$A_p$, $M_p$, $T_p$ & Aggregate residual, pair penalty, exactness threshold \\
$\PR$ & Rosenberg penalty of \eqref{eq:rosenberg} \\
$G_{2q,\sigma}$, $D_{2q,\sigma}$, $S_\sigma$ & Compiled two-qubit gates, two-qubit depth, SWAP count \\
$\mu_i$, $q_{ij}$ & Single and pair pseudo-moments from a common relaxation \\
$L$ & QAOA depth (number of cost and mixer layers) \\
\bottomrule
\end{tabular}
\end{table}

% ---------------------------------------------------------------------
\subsection{Sharp Aggregate-Residual Penalty}
\label{sec:reduction:penalty}

A selective reduction replaces one product $u_iu_j$ by a fresh bit $y_p$ in a chosen
subset of the terms that contain it. The penalty must then dominate the combined effect
of every term sharing $y_p$, not the coefficient of any single term. Let $q(u)=\sum_S d_S u_S$
be a collected multilinear polynomial on $\B^N$, fix a pair $p=\{i,j\}$, and choose a family
$\mathcal{S}_p$ of terms of degree at least three that contain $p$. Writing $v$ for the
coordinates other than $u_i,u_j$, the selected residual is
$A_p(v)=\sum_{S\in\mathcal{S}_p}d_S u_{S\setminus p}$, so that
\begin{equation}
q(u)=B_p(u)+u_iu_jA_p(v),
\label{eq:split}
\end{equation}
where $B_p$ collects all unselected terms. The substituted polynomial is
\begin{equation}
Q_{p,M_p}(u,y_p)=B_p(u)+y_pA_p(v)+M_p\,\PR(u_i,u_j,y_p).
\label{eq:substituted}
\end{equation}

\begin{proposition}[Sharp aggregate-residual penalty]
\label{prop:sharp}
Let $\norminf{A_p}=\max_{v\in\B^{N-2}}|A_p(v)|$ and $M_p\ge 0$. Then
\begin{equation}
q(u)=\min_{y_p\in\B}Q_{p,M_p}(u,y_p)\ \text{for every }u
\iff M_p\ge\norminf{A_p}.
\label{eq:threshold}
\end{equation}
The consistent value $y_p=u_iu_j$ is the unique minimiser for every $u$ if and only if
$M_p>\norminf{A_p}$. Moreover, for every $u$ and $y_p$,
\begin{equation}
Q_{p,M_p}(u,y_p)-q(u)\ \ge\ \bigl(M_p-\norminf{A_p}\bigr)\,\mathbb{1}\{y_p\neq u_iu_j\}.
\label{eq:margin}
\end{equation}
\end{proposition}

The proof evaluates the excess of the unique inconsistent choice of $y_p$ in the three cases
$(u_i,u_j)\in\{(0,0),(0,1)\text{ or }(1,0),(1,1)\}$, which equals $A_p+3M_p$, $A_p+M_p$, and
$M_p-A_p$ respectively. Nonnegativity of all three for every $v$ is equivalent to
\eqref{eq:threshold}, and the bound \eqref{eq:margin} follows from $|A_p|\le\norminf{A_p}$.
Sharpness holds because the maximum defining $\norminf{A_p}$ is attained on the finite cube.
Full details are given in Supplementary Section~S1.

\begin{remark}
Proposition~\ref{prop:sharp} is a sharpened form of the classical Rosenberg condition.
Boros and Gruber \cite{BorosGruber2014} write $f=xyA+B$ and require $M>\max|A|$ for a
substitution applied to all occurrences of $xy$, and they describe an aggregative
approach that groups terms by common parts. The present statement adds three elements
used later in the paper. First, equality $M_p=\norminf{A_p}$ suffices for pointwise
exactness, whereas strict inequality is needed for a unique consistent fibre. Second,
the residual ranges over an arbitrary selected subset $\mathcal{S}_p$, so some
occurrences of $u_iu_j$ may remain native. Third, the separation margin \eqref{eq:margin}
is what makes the compositional argument of the next subsection telescope.
\end{remark}

% ---------------------------------------------------------------------
\subsection{Sequential Compositional Exactness}
\label{sec:reduction:composition}

Successive substitutions may act on terms that already contain earlier auxiliaries, so
exactness must be shown to survive composition. Set $g^{(0)}=f$. At step $r$ choose a pair
of current coordinates, compute the residual $A_r$ from the current collected polynomial
$g^{(r-1)}$ with all coordinates treated as independent Boolean variables, and form
$g^{(r)}$ by \eqref{eq:substituted} with a fresh $y_r$ and penalty $M_r\ge\norminf{A_r}$.
Let $\Phi(x)=(\phi_1(x),\dots,\phi_R(x))$ denote the recursively consistent lift, in which
each $\phi_r(x)$ is the product of the two coordinates replaced at step $r$, evaluated at
$(x,\phi_{<r}(x))$.

\begin{theorem}[Sequential compositional exactness]
\label{thm:composition}
After $R$ substitutions with margins $\delta_r=M_r-\norminf{A_r}\ge 0$,
\begin{equation}
g^{(R)}(x,y)-f(x)\ \ge\ \sum_{r=1}^{R}\delta_r\,\mathbb{1}\{y_r\neq\phi_r(x,y_{<r})\}\ \ge 0
\label{eq:telescope}
\end{equation}
for every $(x,y)$, and $g^{(R)}(x,\Phi(x))=f(x)$. Hence
$f(x)=\min_{y\in\B^R}g^{(R)}(x,y)$ for every $x$. If every $\delta_r>0$, then
$\Phi(x)$ is the unique minimiser.
\end{theorem}

The bound \eqref{eq:margin} applied at each step gives a one-step inequality, and summing
over $r$ telescopes to \eqref{eq:telescope}. Along the consistent lift every penalty
vanishes and every substituted block equals its predecessor, which yields the equality.
Uniqueness follows by considering the first index at which $y$ departs from $\Phi(x)$.
The proof is given in Supplementary Section~S1.

Theorem~\ref{thm:composition} implies that any selected term of degree $k$ can be reduced
to a target locality $\ell\ge 2$ with at most $\max\{k-\ell,0\}$ fresh auxiliaries, while
every unselected term remains native. A chain construction that achieves this bound with
penalties $M\ge|c_A|$ at each step is given in Supplementary Section~S1. Sharing an
intermediate product between terms can lower the auxiliary count, but the shared penalty
must then be recomputed from the aggregate residual. This is the standard recursive
Rosenberg construction \cite{Rosenberg1975,BorosGruber2014}, and the contribution here is
the explicit selective form with a pointwise guarantee at each stage.

% ---------------------------------------------------------------------
\subsection{Cubic Closed Form and Pair Reuse}
\label{sec:reduction:cubic}

Cubic objectives make the threshold explicit and are the setting of the experiments.
Write $f=f_{\le 2}+\sum_{A\in\mathcal{C}_3(f)}c_Ax_A$, where
$\mathcal{C}_3(f)=\{A\subseteq[n] : |A|=3,\ c_A\neq 0\}$. A design $\sigma=(D,\rho)$
chooses the reduced supports $D\subseteq\mathcal{C}_3(f)$ and, for each $A\in D$, one pair
$\rho(A)\in\binom{A}{2}$. Let $R=\mathcal{C}_3(f)\setminus D$, let $P_\rho=\rho(D)$ be the
set of active pairs, and let $D_p=\{A\in D : \rho(A)=p\}$. With
$A_p(x)=\sum_{A\in D_p}c_Ax_{A\setminus p}$ the representation is
\begin{equation}
F_{\rho,M}(x,y)=f_{\le 2}(x)+\!\sum_{A\in R}c_Ax_A
+\!\sum_{p\in P_\rho}\!\bigl[y_pA_p(x)+M_p\PR(x_i,x_j,y_p)\bigr].
\label{eq:cubic_rep}
\end{equation}
Each $A_p$ is linear in distinct remaining coordinates, so its Boolean sup norm is the
larger of the total positive and total negative coefficient mass,
\begin{equation}
T_p=\max\Bigl\{\sum_{A\in D_p,\,c_A>0}c_A,\ \sum_{A\in D_p,\,c_A<0}|c_A|\Bigr\}=\norminf{A_p}.
\label{eq:cubic_threshold}
\end{equation}

\begin{proposition}[Exact cubic mixed-sign reduction]
\label{prop:cubic}
For nonnegative penalties, $f(x)=\min_yF_{\rho,M}(x,y)$ for every $x$ if and only if
$M_p\ge T_p$ for every active pair. Strict inequalities give the unique consistent lift
$\phi_\rho(x)=(x_ix_j)_{p=\{i,j\}\in P_\rho}$.
\end{proposition}

Minimisation separates over the auxiliaries, so the claim follows from
Proposition~\ref{prop:sharp} applied to each pair. Fixing positive margins
$M_p=T_p+\varepsilon_p$ defines the finite exact family $\Sigma^{\varepsilon}_{\mathrm{ex}}$
used in the experiments. Its endpoints are the native design $D=\emptyset$ and the fully
reduced designs $D=\mathcal{C}_3(f)$. The auxiliary count is $a_\sigma=|P_\rho|$ rather than
$|D|$, because several cubic terms sharing a pair use one auxiliary. The minimum number of
auxiliaries that reduces all of $D$ is the set-cover number of the pair shadow
$\partial_2D=\bigcup_{A\in D}\binom{A}{2}$, which satisfies
$\lceil |D|/\Delta_2(D)\rceil\le\tau_2(D)\le|D|$ with $\Delta_2(D)$ the maximum number of
terms containing a common pair \cite{VermaLewisKochenberger2021,Schmidbauer2026}.
Supplementary Section~S1 gives the integer-programming form and the overlap-graph view.

% ---------------------------------------------------------------------
\subsection{Worked Max-3SAT Example}
\label{sec:reduction:example}

A small instance illustrates the mixed-sign threshold and the three representations that the
experiments compare. Consider the weighted clauses
$C_1=(\neg x_1\vee\neg x_2\vee\neg x_3)$, $C_2=(x_2\vee x_3\vee x_4)$, and
$C_3=(\neg x_1\vee\neg x_4\vee\neg x_5)$ with weights $5$, $4$, and $2$. The total violated
weight expands and aggregates to
\begin{equation}
f^{\mathrm{ex}}(x)=g(x)+5x_1x_2x_3-4x_2x_3x_4+2x_1x_4x_5,
\label{eq:example_f}
\end{equation}
where $g(x)=4-4(x_2+x_3+x_4)+4(x_2x_3+x_2x_4+x_3x_4)$ is quadratic. The first two cubic
terms share the pair $\{2,3\}$ with opposite signs. Substituting $y_{23}=x_2x_3$ gives the
aggregate residual $A_{23}=5x_1-4x_4$, so \eqref{eq:cubic_threshold} yields
$T_{23}=\max\{5,4\}=5$. Choosing $M_{23}=6$ produces the selective representation
\begin{equation}
F_{\mathrm{sel}}=g(x)+5x_1y_{23}-4y_{23}x_4+2x_1x_4x_5+6\PR(x_2,x_3,y_{23}),
\label{eq:example_sel}
\end{equation}
which retains one cubic term natively. Full quadratization additionally substitutes
$y_{14}=x_1x_4$ in the remaining term with $T_{14}=2$ and $M_{14}=3$. Exhaustive enumeration
over all $32$ original assignments confirms that all three representations satisfy
\eqref{eq:pointwise} with unique consistent minimisers.

Table~\ref{tab:example} compares the logical resources. The selective choice removes two
cubic terms with one auxiliary and saves one auxiliary and four quadratic couplings relative
to full reduction. Fig.~\ref{fig:example_circuits} shows one cost layer of each representation
under the parity-gadget synthesis of Section~\ref{sec:selection:model}. These logical and
reference-circuit counts do not by themselves establish a compiled or variational advantage,
which is the question addressed in Section~\ref{sec:experiments}.

\begin{table}[!t]
\caption{Logical comparison for the worked Max-3SAT instance. Couplings count distinct
degree-two monomials after aggregation.}
\label{tab:example}
\centering
\footnotesize
\begin{tabular}{lccccc}
\toprule
Encoding & Qubits & Aux. & Cubic terms & Couplings & Penalties \\
\midrule
Native & 5 & 0 & 3 & 3 & -- \\
Selective & 6 & 1 & 1 & 7 & $M_{23}=6$ \\
Fully reduced & 7 & 2 & 0 & 11 & $M_{23}=6$, $M_{14}=3$ \\
\bottomrule
\end{tabular}
\end{table}

\begin{figure}[!t]
\centering
% \includegraphics[width=\columnwidth]{fig_example_circuits.pdf}
\fbox{\parbox{0.95\columnwidth}{\centering\vspace{2.2cm}Fig.~2 placeholder: one cost layer
for the native, selective, and fully reduced encodings of \eqref{eq:example_f}, drawn with
the parity-gadget synthesis, with CNOT counts annotated.\vspace{2.2cm}}}
\caption{One QAOA cost layer for the three representations of the worked instance in
\eqref{eq:example_f}. Native synthesis uses three four-CNOT cubic gadgets. The selective
encoding replaces two of them by an auxiliary qubit with two-qubit penalty phases and
retains one cubic gadget. The fully reduced encoding uses two auxiliaries and no cubic gadget.}
\label{fig:example_circuits}
\end{figure}

% =====================================================================
\section{Resource-Aware Representation Selection}
\label{sec:selection}

\subsection{Compilation and Circuit Model}
\label{sec:selection:model}

\subsection{Resource Score, Budgets, and Fibre Diagnostic}
\label{sec:selection:score}

\subsection{Pareto Beam Search}
\label{sec:selection:beam}

\subsection{Branch-and-Bound Certification}
\label{sec:selection:certification}

\subsection{Representation-Consistent Warm Starts}
\label{sec:selection:warmstart}

% Fig. 3: selection pipeline (beam, optional certification, warm-start output).

% =====================================================================
\section{Experimental Evaluation}
\label{sec:experiments}

\subsection{Benchmark Families and Instance Tiers}
\label{sec:experiments:benchmarks}

\subsection{Representations, Baselines, and Compilation Protocol}
\label{sec:experiments:protocol}

\subsection{QAOA Protocol and Statistical Treatment}
\label{sec:experiments:qaoa_protocol}

\subsection{Exactness and Selector Validation}
\label{sec:experiments:validation}

\subsection{Compiled Resources Across Coupling Maps}
\label{sec:experiments:resources}

\subsection{Finite-Depth QAOA Under Matched Budgets}
\label{sec:experiments:qaoa}

\subsection{Warm-Start Performance}
\label{sec:experiments:warmstart}

% Fig. 4: resource results (paired differences vs. native; selective vs. random on five maps).
% Fig. 5: paired QAOA objective differences, positive depth only, with inconsistency inset.
% Fig. 6: warm-minus-cold on the screened set with LP-rounding reference.

% =====================================================================
\section{Discussion}
\label{sec:discussion}

% =====================================================================
\section{Conclusion}
\label{sec:conclusion}

% =====================================================================
\section*{Acknowledgment}
The authors thank the University of Warwick Undergraduate Research Support Scheme
for supporting this project.

\section*{Data and Code Availability}
Processed data, source manifests, and the code used to generate and analyze the
reported results are available in the project repository. The complete raw result
archive will be deposited with a persistent identifier before publication.

% =====================================================================
% =====================================================================
% bibliography.tex  --  reference list for main.tex (input by % =====================================================================
% bibliography.tex  --  reference list for main.tex (input by % =====================================================================
% bibliography.tex  --  reference list for main.tex (input by \input{bibliography})
%
% The entries below were carried over from the original manuscript
% (quantumarticle format). IEEE TQE uses numbered references ordered by
% first citation. Because the list is typeset with thebibliography rather
% than BibTeX, the \bibitem order must match citation order in the final
% text and each entry must be reformatted to IEEE reference style, e.g.
%   E. Boros and P. L. Hammer, "Pseudo-Boolean optimization," Discrete Appl.
%   Math., vol. 123, no. 1-3, pp. 155-225, 2002.
% Newly added references requested in review appear at the end of the list
% under "Added in revision". All new entries must be verified against the
% published versions before submission.
% =====================================================================
\begin{thebibliography}{99}

\bibitem{BorosHammer2002}
Endre Boros and Peter~L. Hammer.
\newblock ``Pseudo-boolean optimization''.
\newblock Discrete Applied Mathematics {\bf 123}, 155--225~(2002).

\bibitem{Hadfield2021}
Stuart Hadfield.
\newblock ``On the representation of Boolean and real functions as Hamiltonians for quantum computing''.
\newblock ACM Transactions on Quantum Computing {\bf 2}, Article 18, 1--21~(2021).
\newblock arXiv:1804.09130.

\bibitem{Farhi2014}
Edward Farhi, Jeffrey Goldstone, and Sam Gutmann.
\newblock ``A quantum approximate optimization algorithm''~(2014).
\newblock arXiv:1411.4028.

\bibitem{Zhou2020}
Leo Zhou, Sheng-Tao Wang, Soonwon Choi, Hannes Pichler, and Mikhail D. Lukin.
\newblock ``Quantum Approximate Optimization Algorithm: Performance, Mechanism, and Implementation on Near-Term Devices''.
\newblock Physical Review X {\bf 10}, 021067~(2020).


\bibitem{Glos2022}
Adam Glos, Aleksandra Krawiec, and Zolt{\'a}n Zimbor{\'a}s.
\newblock ``Space-efficient binary optimization for variational quantum computing''.
\newblock npj Quantum Information {\bf 8}, 39~(2022).
\newblock arXiv:2009.07309.

\bibitem{CampbellDahl2022}
Colin Campbell and Edward Dahl.
\newblock ``QAOA of the highest order''.
\newblock In 2022 IEEE 19th International Conference on Software Architecture Companion (ICSA-C), pages 141--146~(2022).
\newblock doi:10.1109/ICSA-C54293.2022.00035.
\newblock arXiv:2111.12754.

\bibitem{Stein2023}
Jonas Stein, Farbod Chamanian, Maximilian Zorn, Jonas N{\"u}{\ss}lein, Sebastian Zielinski, Michael K{\"o}lle, and Claudia Linnhoff-Popien.
\newblock ``Evidence that PUBO outperforms QUBO when solving continuous optimization problems with the QAOA''.
\newblock In Proceedings of the Genetic and Evolutionary Computation Conference Companion.
\newblock Pages 2254--2262~(2023).
\newblock arXiv:2305.03390.

\bibitem{Bell2025}
Kristina Bell, Adam Lowe, Emily Coles, Nathanael Ridgway, and Gillian Marshall.
\newblock ``A comparison of quadratic and higher-order representations for QAOA''~(2025).
\newblock arXiv:2509.20127.

\bibitem{Dragoi2026}
Sabina Dr{\u a}goi, Alberto Baiardi, and Daniel~J. Egger.
\newblock ``Approximate quadratization of high-order Hamiltonians for combinatorial quantum optimization''.
\newblock Physical Review Research {\bf 8}, 023159~(2026).
\newblock arXiv:2505.04700.

\bibitem{Verchere2023}
Zo{\'e} Verch{\`e}re, Sourour Elloumi, and Andrea Simonetto.
\newblock ``Optimizing variational circuits for higher-order binary optimization''.
\newblock In 2023 IEEE International Conference on Quantum Computing and Engineering (QCE), volume 1, pages 19--25~(2023).
\newblock doi:10.1109/QCE57702.2023.00011.
\newblock arXiv:2307.16756.

\bibitem{Weidenfeller2022}
Johannes Weidenfeller, Lucia~C. Valor, Julien Gacon, Caroline Tornow, Luciano Bello, Stefan Woerner, and Daniel~J. Egger.
\newblock ``Scaling of the quantum approximate optimization algorithm on superconducting qubit based hardware''.
\newblock Quantum {\bf 6}, 870~(2022).
\newblock arXiv:2202.03459.

\bibitem{Harrigan2021}
Matthew P. Harrigan et al.
\newblock ``Quantum approximate optimization of non-planar graph problems on a planar superconducting processor''.
\newblock Nature Physics {\bf 17}, 332--336~(2021).


\bibitem{Rosenberg1975}
I.~G. Rosenberg.
\newblock ``Reduction of bivalent maximization to the quadratic case''.
\newblock Cahiers du Centre d'\'Etudes de Recherche Op\'erationnelle {\bf 17}, 71--74~(1975).

\bibitem{BorosGruber2014}
Endre Boros and Aritanan Gruber.
\newblock ``On quadratization of pseudo-boolean functions''~(2014).
\newblock arXiv:1404.6538.

\bibitem{Dattani2019}
Nike Dattani.
\newblock ``Quadratization in discrete optimization and quantum mechanics''~(2019).
\newblock arXiv:1901.04405.

\bibitem{FreedmanDrineas2005}
Daniel Freedman and Petros Drineas.
\newblock ``Energy minimization via graph cuts: Settling what is possible''.
\newblock In 2005 IEEE Computer Society Conference on Computer Vision and Pattern Recognition (CVPR'05), volume 2, pages 939--946~(2005).
\newblock doi:10.1109/CVPR.2005.143.

\bibitem{Ishikawa2011}
Hiroshi Ishikawa.
\newblock ``Transformation of general binary MRF minimization to the first-order case''.
\newblock IEEE Transactions on Pattern Analysis and Machine Intelligence {\bf 33}, 1234--1249~(2011).

\bibitem{AnthonyEtAl2017}
Martin Anthony, Endre Boros, Yves Crama, and Aritanan Gruber.
\newblock ``Quadratic reformulations of nonlinear binary optimization problems''.
\newblock Mathematical Programming {\bf 162}, 115--144~(2017).

\bibitem{Anthony2016}
Martin Anthony, Endre Boros, Yves Crama, and Aritanan Gruber.
\newblock ``Quadratization of symmetric pseudo-Boolean functions''.
\newblock Discrete Applied Mathematics {\bf 203}, 1--12~(2016).


\bibitem{BorosCramaRodriguezHeck2020}
Endre Boros, Yves Crama, and Elisabeth Rodr{\'i}guez-Heck.
\newblock ``Compact quadratizations for pseudo-Boolean functions''.
\newblock Journal of Combinatorial Optimization {\bf 39}, 687--707~(2020).

\bibitem{VermaLewis2020}
Amit Verma and Mark Lewis.
\newblock ``Optimal quadratic reformulations of fourth degree pseudo-Boolean functions''.
\newblock Optimization Letters {\bf 14}, 1557--1569~(2020).

\bibitem{VermaLewisKochenberger2021}
Amit Verma, Mark Lewis, and Gary Kochenberger.
\newblock ``Efficient QUBO transformation for higher degree pseudo Boolean functions''~(2021).
\newblock arXiv:2107.11695.

\bibitem{CramaEtAl2025}
Yves Crama, Sourour Elloumi, Am{\'e}lie Lambert, and Elisabeth Rodr{\'i}guez-Heck.
\newblock ``Quadratization and convexification in polynomial binary optimization''.
\newblock Journal of Combinatorial Optimization {\bf 50}, Article 28~(2025).

\bibitem{Herrman2021}
Rebekah Herrman, Lorna Treffert, James Ostrowski, Phillip~C. Lotshaw, Travis~S. Humble, and George Siopsis.
\newblock ``Globally optimizing QAOA circuit depth for constrained optimization problems''.
\newblock Algorithms {\bf 14}, 294~(2021).
\newblock arXiv:2108.03281.

\bibitem{Schmidbauer2024}
Lukas Schmidbauer, Karen Wintersperger, Elisabeth Lobe, and Wolfgang Mauerer.
\newblock ``Polynomial reduction methods and their impact on QAOA circuits''.
\newblock In 2024 IEEE International Conference on Quantum Software (QSW), pages 35--45~(2024).
\newblock doi:10.1109/QSW62656.2024.00018.
\newblock arXiv:2406.08889.

\bibitem{Schmidbauer2026}
Lukas Schmidbauer, Elisabeth Lobe, Ina Schaefer, and Wolfgang Mauerer.
\newblock ``It's quick to be square: Fast quadratisation for quantum toolchains''.
\newblock ACM Transactions on Quantum Computing {\bf 7}, Article 18, 1--46~(2026).
\newblock arXiv:2411.19934.

\bibitem{Rovara2026}
Damian Rovara, Lukas Burgholzer, and Robert Wille.
\newblock ``Quantum hardware-efficient selection of auxiliary variables for QUBO formulations''.
\newblock In 2026 Design, Automation \& Test in Europe Conference (DATE), pages 1--7~(2026).
\newblock doi:10.23919/DATE69613.2026.11539407.
\newblock arXiv:2511.19613.

\bibitem{Egger2021}
Daniel~J. Egger, Jakub Mare{\v c}ek, and Stefan Woerner.
\newblock ``Warm-starting quantum optimization''.
\newblock Quantum {\bf 5}, 479~(2021).

\bibitem{TateSDP2023}
Reuben Tate, Majid Farhadi, Creston Herold, Greg Mohler, and Swati Gupta.
\newblock ``Bridging classical and quantum with SDP initialized warm-starts for QAOA''.
\newblock ACM Transactions on Quantum Computing {\bf 4}, Article 9, 1--39~(2023).

\bibitem{TateCustom2023}
Reuben Tate, Jai Moondra, Bryan Gard, Greg Mohler, and Swati Gupta.
\newblock ``Warm-started QAOA with custom mixers provably converges and computationally beats Goemans--Williamson's Max-Cut at low circuit depths''.
\newblock Quantum {\bf 7}, 1121~(2023).

\bibitem{WangMaxCut2018}
Zhihui Wang, Stuart Hadfield, Zhang Jiang, and Eleanor G. Rieffel.
\newblock ``Quantum approximate optimization algorithm for MaxCut: A fermionic view''.
\newblock Physical Review A {\bf 97}, 022304~(2018).


\bibitem{FarhiLocal2020}
Edward Farhi, David Gamarnik, and Sam Gutmann.
\newblock ``The Quantum Approximate Optimization Algorithm Needs to See the Whole Graph: Worst Case Examples''.
\newblock Preprint at arXiv:2005.08747~(2020).


\bibitem{Akshay2020}
V. Akshay, H. Philathong, M. E. S. Morales, and J. D. Biamonte.
\newblock ``Reachability Deficits in Quantum Approximate Optimization''.
\newblock Physical Review Letters {\bf 124}, 090504~(2020).


\bibitem{Cain2023}
Madelyn Cain, Edward Farhi, Sam Gutmann, Daniel Ranard, and Eugene Tang.
\newblock ``The QAOA gets stuck starting from a good classical string''.
\newblock Preprint at arXiv:2207.05089v2~(2023).


\bibitem{LandDoig1960}
A. H. Land and A. G. Doig.
\newblock ``An Automatic Method of Solving Discrete Programming Problems''.
\newblock Econometrica {\bf 28}, 497--520~(1960).


\bibitem{LawlerWood1966}
E. L. Lawler and D. E. Wood.
\newblock ``Branch-and-Bound Methods: A Survey''.
\newblock Operations Research {\bf 14}, 699--719~(1966).


\bibitem{Hadfield2019}
Stuart Hadfield, Zhihui Wang, Bryan O'Gorman, Eleanor G. Rieffel, Davide Venturelli, and Rupak Biswas.
\newblock ``From the Quantum Approximate Optimization Algorithm to a Quantum Alternating Operator Ansatz''.
\newblock Algorithms {\bf 12}, 34~(2019).


\bibitem{Powell1994}
M. J. D. Powell.
\newblock ``A Direct Search Optimization Method That Models the Objective and Constraint Functions by Linear Interpolation''.
\newblock In Advances in Optimization and Numerical Analysis, pages 51--67~(1994).


\bibitem{Student1908}
Student.
\newblock ``The Probable Error of a Mean''.
\newblock Biometrika {\bf 6}, 1--25~(1908).


\bibitem{Lucas2014}
Andrew Lucas.
\newblock ``Ising formulations of many NP problems''.
\newblock Frontiers in Physics {\bf 2}, 5~(2014).


\bibitem{Glover2018}
Fred Glover, Gary Kochenberger, and Yu Du.
\newblock ``A Tutorial on Formulating and Using QUBO Models''.
\newblock Preprint at arXiv:1811.11538~(2018).


\bibitem{CowtanPhase2020}
Alexander Cowtan, Silas Dilkes, Ross Duncan, Will Simmons, and Seyon Sivarajah.
\newblock ``Phase Gadget Synthesis for Shallow Circuits''.
\newblock Electronic Proceedings in Theoretical Computer Science {\bf 318}, 213--228~(2020).


\bibitem{CowtanUCC2020}
Alexander Cowtan, Will Simmons, and Ross Duncan.
\newblock ``A Generic Compilation Strategy for the Unitary Coupled Cluster Ansatz''.
\newblock Preprint at arXiv:2007.10515~(2020).


\bibitem{Murali2019}
Prakash Murali, Jonathan M. Baker, Ali Javadi-Abhari, Frederic T. Chong, and Margaret Martonosi.
\newblock ``Noise-Adaptive Compiler Mappings for Noisy Intermediate-Scale Quantum Computers''.
\newblock Proceedings of the Twenty-Fourth International Conference on Architectural Support for Programming Languages and Operating Systems, 1015--1029~(2019).


\bibitem{DebEtAl2002}
Kalyanmoy Deb, Amrit Pratap, Sameer Agarwal, and T.~Meyarivan.
\newblock ``A fast and elitist multiobjective genetic algorithm: NSGA-II''.
\newblock IEEE Transactions on Evolutionary Computation {\bf 6}, 182--197~(2002).

\bibitem{Basso2022}
Joao Basso, David Gamarnik, Song Mei, and Leo Zhou.
\newblock ``Performance and limitations of the QAOA at constant levels on large sparse hypergraphs and spin glass models''.
\newblock In 2022 IEEE 63rd Annual Symposium on Foundations of Computer Science (FOCS), pages 335--343~(2022).
\newblock doi:10.1109/FOCS54457.2022.00039.
\newblock arXiv:2204.10306.

\bibitem{PelofskeShortDepth2024}
Elijah Pelofske, Andreas B{\"a}rtschi, and Stephan Eidenbenz.
\newblock ``Short-depth QAOA circuits and quantum annealing on higher-order Ising models''.
\newblock npj Quantum Information {\bf 10}, 30~(2024).

\bibitem{Pelofske2024}
Elijah Pelofske, Andreas B{\"a}rtschi, Lukasz Cincio, John Golden, and Stephan Eidenbenz.
\newblock ``Scaling whole-chip QAOA for higher-order Ising spin glass models on heavy-hex graphs''.
\newblock npj Quantum Information {\bf 10}, 109~(2024).
\newblock arXiv:2312.00997.

\bibitem{Qiskit2024}
Ali Javadi-Abhari et al.
\newblock ``Quantum computing with Qiskit''.
\newblock Preprint at arXiv:2405.08810~(2024).


\bibitem{Virtanen2020}
Pauli Virtanen et al.
\newblock ``SciPy 1.0: fundamental algorithms for scientific computing in Python''.
\newblock Nature Methods {\bf 17}, 261--272~(2020).


\bibitem{Li2019}
Gushu Li, Yufei Ding, and Yuan Xie.
\newblock ``Tackling the Qubit Mapping Problem for NISQ-Era Quantum Devices''.
\newblock Proceedings of the Twenty-Fourth International Conference on Architectural Support for Programming Languages and Operating Systems, 1001--1014~(2019).


\bibitem{Crooks2018}
Gavin E. Crooks.
\newblock ``Performance of the Quantum Approximate Optimization Algorithm on the Maximum Cut Problem''.
\newblock Preprint at arXiv:1811.08419~(2018).


\bibitem{Sack2021}
Stefan H. Sack and Maksym Serbyn.
\newblock ``Quantum annealing initialization of the quantum approximate optimization algorithm''.
\newblock Quantum {\bf 5}, 491~(2021).


\bibitem{HuangfuHall2018}
Q. Huangfu and J. A. J. Hall.
\newblock ``Parallelizing the dual revised simplex method''.
\newblock Mathematical Programming Computation {\bf 10}, 119--142~(2018).


\bibitem{Marshall2020}
Jeffrey Marshall, Filip Wudarski, Stuart Hadfield, and Tad Hogg.
\newblock ``Characterizing local noise in QAOA circuits''.
\newblock IOP SciNotes {\bf 1}, 025208~(2020).


\bibitem{WangNoise2021}
Samson Wang et al.
\newblock ``Noise-induced barren plateaus in variational quantum algorithms''.
\newblock Nature Communications {\bf 12}, 6961~(2021).


\bibitem{Franca2021}
Daniel Stilck França and Raul García-Patrón.
\newblock ``Limitations of optimization algorithms on noisy quantum devices''.
\newblock Nature Physics {\bf 17}, 1221--1227~(2021).


\bibitem{Nagies2025}
Sebastian Nagies, Kevin~T. Geier, Javed Akram, Dimitrios Bantounas, Michael Johanning, and Philipp Hauke.
\newblock ``Boosting quantum annealing performance through direct polynomial unconstrained binary optimization''.
\newblock Quantum Science and Technology {\bf 10}, 035008~(2025).
\newblock arXiv:2412.04398.

\bibitem{Nusslein2023}
Jonas N{\"u}{\ss}lein, Sebastian Zielinski, Thomas Gabor, Claudia Linnhoff-Popien, and Sebastian Feld.
\newblock ``Solving (Max) 3-SAT via quadratic unconstrained binary optimization''~(2023).
\newblock arXiv:2302.03536.

\bibitem{Wurtz2021}
Jonathan Wurtz and Peter Love.
\newblock ``MaxCut quantum approximate optimization algorithm performance guarantees for $p>1$''.
\newblock Physical Review A {\bf 103}, 042612~(2021).


\bibitem{BassoHighDepth2022}
Joao Basso, Edward Farhi, Kunal Marwaha, Benjamin Villalonga, and Leo Zhou.
\newblock ``The Quantum Approximate Optimization Algorithm at High Depth for MaxCut on Large-Girth Regular Graphs and the Sherrington-Kirkpatrick Model''.
\newblock Leibniz International Proceedings in Informatics (TQC 2022) {\bf 232}, 7:1--7:21~(2022).


\bibitem{SheraliAdams1990}
Hanif D. Sherali and Warren P. Adams.
\newblock ``A Hierarchy of Relaxations between the Continuous and Convex Hull Representations for Zero-One Programming Problems''.
\newblock SIAM Journal on Discrete Mathematics {\bf 3}, 411--430~(1990).


\bibitem{McCormick1976}
Garth P. McCormick.
\newblock ``Computability of global solutions to factorable nonconvex programs: Part I --- Convex underestimating problems''.
\newblock Mathematical Programming {\bf 10}, 147--175~(1976).


% -------------------- Added in revision --------------------

\bibitem{Koch2026}
Frederik Koch, Shahram Panahiyan, Rick Mukherjee, Joseph Doetsch, et al.
\newblock ``Resource-efficient quantum optimization via higher-order encoding''.
\newblock EPJ Quantum Technology~(2026).
\newblock arXiv:2511.17545.

\bibitem{Mirkarimi2024}
Puya Mirkarimi, David C. Hoyle, Ross Williams, and Nicholas Chancellor.
\newblock ``Quantum optimization with linear Ising penalty functions for customer data science''.
\newblock New Journal of Physics {\bf 26}, 103005~(2024).

\bibitem{Sriluckshmy2023}
P.~V. Sriluckshmy, Vicente Pina-Canelles, Mario Ponce, Manuel~G. Algaba, Fedor {\v S}imkovic, and Martin Leib.
\newblock ``Optimal, hardware native decomposition of parameterized multi-qubit Pauli gates''.
\newblock Quantum Science and Technology {\bf 8}, 045029~(2023).

\bibitem{Amy2018}
Matthew Amy, Parsiad Azimzadeh, and Michele Mosca.
\newblock ``On the CNOT-complexity of CNOT-PHASE circuits''.
\newblock Quantum Science and Technology {\bf 4}, 015002~(2018).

\bibitem{Holmes2020}
Adam Holmes, Sonika Johri, Gian~Giacomo Guerreschi, James~S. Clarke, and A.~Y. Matsuura.
\newblock ``Impact of qubit connectivity on quantum algorithm performance''.
\newblock Quantum Science and Technology {\bf 5}, 025009~(2020).

\bibitem{Zulehner2019}
Alwin Zulehner, Alexandru Paler, and Robert Wille.
\newblock ``An efficient methodology for mapping quantum circuits to the IBM QX architectures''.
\newblock IEEE Transactions on Computer-Aided Design of Integrated Circuits and Systems {\bf 38}, 1226--1236~(2019).

\bibitem{Tan2021}
Benjamin Tan, Marc-Antoine Lemonde, Supanut Thanasilp, Jirawat Tangpanitanon, and Dimitris~G. Angelakis.
\newblock ``Qubit-efficient encoding schemes for binary optimisation problems''.
\newblock Quantum {\bf 5}, 454~(2021).

\bibitem{Herrman2021depth}
Rebekah Herrman, James Ostrowski, Travis~S. Humble, and George Siopsis.
\newblock ``Lower bounds on circuit depth of the quantum approximate optimization algorithm''.
\newblock Quantum Information Processing {\bf 20}, 59~(2021).

\bibitem{Abbas2024}
Amira Abbas et al.
\newblock ``Challenges and opportunities in quantum optimization''.
\newblock Nature Reviews Physics {\bf 6}, 718--735~(2024).

\bibitem{Ayodele2022}
Mayowa Ayodele.
\newblock ``Penalty weights in QUBO formulations: Permutation problems''.
\newblock In Evolutionary Computation in Combinatorial Optimization (EvoCOP 2022), Lecture Notes in Computer Science {\bf 13222}, pages 159--174~(2022).

\end{thebibliography}
)
%
% The entries below were carried over from the original manuscript
% (quantumarticle format). IEEE TQE uses numbered references ordered by
% first citation. Because the list is typeset with thebibliography rather
% than BibTeX, the \bibitem order must match citation order in the final
% text and each entry must be reformatted to IEEE reference style, e.g.
%   E. Boros and P. L. Hammer, "Pseudo-Boolean optimization," Discrete Appl.
%   Math., vol. 123, no. 1-3, pp. 155-225, 2002.
% Newly added references requested in review appear at the end of the list
% under "Added in revision". All new entries must be verified against the
% published versions before submission.
% =====================================================================
\begin{thebibliography}{99}

\bibitem{BorosHammer2002}
Endre Boros and Peter~L. Hammer.
\newblock ``Pseudo-boolean optimization''.
\newblock Discrete Applied Mathematics {\bf 123}, 155--225~(2002).

\bibitem{Hadfield2021}
Stuart Hadfield.
\newblock ``On the representation of Boolean and real functions as Hamiltonians for quantum computing''.
\newblock ACM Transactions on Quantum Computing {\bf 2}, Article 18, 1--21~(2021).
\newblock arXiv:1804.09130.

\bibitem{Farhi2014}
Edward Farhi, Jeffrey Goldstone, and Sam Gutmann.
\newblock ``A quantum approximate optimization algorithm''~(2014).
\newblock arXiv:1411.4028.

\bibitem{Zhou2020}
Leo Zhou, Sheng-Tao Wang, Soonwon Choi, Hannes Pichler, and Mikhail D. Lukin.
\newblock ``Quantum Approximate Optimization Algorithm: Performance, Mechanism, and Implementation on Near-Term Devices''.
\newblock Physical Review X {\bf 10}, 021067~(2020).


\bibitem{Glos2022}
Adam Glos, Aleksandra Krawiec, and Zolt{\'a}n Zimbor{\'a}s.
\newblock ``Space-efficient binary optimization for variational quantum computing''.
\newblock npj Quantum Information {\bf 8}, 39~(2022).
\newblock arXiv:2009.07309.

\bibitem{CampbellDahl2022}
Colin Campbell and Edward Dahl.
\newblock ``QAOA of the highest order''.
\newblock In 2022 IEEE 19th International Conference on Software Architecture Companion (ICSA-C), pages 141--146~(2022).
\newblock doi:10.1109/ICSA-C54293.2022.00035.
\newblock arXiv:2111.12754.

\bibitem{Stein2023}
Jonas Stein, Farbod Chamanian, Maximilian Zorn, Jonas N{\"u}{\ss}lein, Sebastian Zielinski, Michael K{\"o}lle, and Claudia Linnhoff-Popien.
\newblock ``Evidence that PUBO outperforms QUBO when solving continuous optimization problems with the QAOA''.
\newblock In Proceedings of the Genetic and Evolutionary Computation Conference Companion.
\newblock Pages 2254--2262~(2023).
\newblock arXiv:2305.03390.

\bibitem{Bell2025}
Kristina Bell, Adam Lowe, Emily Coles, Nathanael Ridgway, and Gillian Marshall.
\newblock ``A comparison of quadratic and higher-order representations for QAOA''~(2025).
\newblock arXiv:2509.20127.

\bibitem{Dragoi2026}
Sabina Dr{\u a}goi, Alberto Baiardi, and Daniel~J. Egger.
\newblock ``Approximate quadratization of high-order Hamiltonians for combinatorial quantum optimization''.
\newblock Physical Review Research {\bf 8}, 023159~(2026).
\newblock arXiv:2505.04700.

\bibitem{Verchere2023}
Zo{\'e} Verch{\`e}re, Sourour Elloumi, and Andrea Simonetto.
\newblock ``Optimizing variational circuits for higher-order binary optimization''.
\newblock In 2023 IEEE International Conference on Quantum Computing and Engineering (QCE), volume 1, pages 19--25~(2023).
\newblock doi:10.1109/QCE57702.2023.00011.
\newblock arXiv:2307.16756.

\bibitem{Weidenfeller2022}
Johannes Weidenfeller, Lucia~C. Valor, Julien Gacon, Caroline Tornow, Luciano Bello, Stefan Woerner, and Daniel~J. Egger.
\newblock ``Scaling of the quantum approximate optimization algorithm on superconducting qubit based hardware''.
\newblock Quantum {\bf 6}, 870~(2022).
\newblock arXiv:2202.03459.

\bibitem{Harrigan2021}
Matthew P. Harrigan et al.
\newblock ``Quantum approximate optimization of non-planar graph problems on a planar superconducting processor''.
\newblock Nature Physics {\bf 17}, 332--336~(2021).


\bibitem{Rosenberg1975}
I.~G. Rosenberg.
\newblock ``Reduction of bivalent maximization to the quadratic case''.
\newblock Cahiers du Centre d'\'Etudes de Recherche Op\'erationnelle {\bf 17}, 71--74~(1975).

\bibitem{BorosGruber2014}
Endre Boros and Aritanan Gruber.
\newblock ``On quadratization of pseudo-boolean functions''~(2014).
\newblock arXiv:1404.6538.

\bibitem{Dattani2019}
Nike Dattani.
\newblock ``Quadratization in discrete optimization and quantum mechanics''~(2019).
\newblock arXiv:1901.04405.

\bibitem{FreedmanDrineas2005}
Daniel Freedman and Petros Drineas.
\newblock ``Energy minimization via graph cuts: Settling what is possible''.
\newblock In 2005 IEEE Computer Society Conference on Computer Vision and Pattern Recognition (CVPR'05), volume 2, pages 939--946~(2005).
\newblock doi:10.1109/CVPR.2005.143.

\bibitem{Ishikawa2011}
Hiroshi Ishikawa.
\newblock ``Transformation of general binary MRF minimization to the first-order case''.
\newblock IEEE Transactions on Pattern Analysis and Machine Intelligence {\bf 33}, 1234--1249~(2011).

\bibitem{AnthonyEtAl2017}
Martin Anthony, Endre Boros, Yves Crama, and Aritanan Gruber.
\newblock ``Quadratic reformulations of nonlinear binary optimization problems''.
\newblock Mathematical Programming {\bf 162}, 115--144~(2017).

\bibitem{Anthony2016}
Martin Anthony, Endre Boros, Yves Crama, and Aritanan Gruber.
\newblock ``Quadratization of symmetric pseudo-Boolean functions''.
\newblock Discrete Applied Mathematics {\bf 203}, 1--12~(2016).


\bibitem{BorosCramaRodriguezHeck2020}
Endre Boros, Yves Crama, and Elisabeth Rodr{\'i}guez-Heck.
\newblock ``Compact quadratizations for pseudo-Boolean functions''.
\newblock Journal of Combinatorial Optimization {\bf 39}, 687--707~(2020).

\bibitem{VermaLewis2020}
Amit Verma and Mark Lewis.
\newblock ``Optimal quadratic reformulations of fourth degree pseudo-Boolean functions''.
\newblock Optimization Letters {\bf 14}, 1557--1569~(2020).

\bibitem{VermaLewisKochenberger2021}
Amit Verma, Mark Lewis, and Gary Kochenberger.
\newblock ``Efficient QUBO transformation for higher degree pseudo Boolean functions''~(2021).
\newblock arXiv:2107.11695.

\bibitem{CramaEtAl2025}
Yves Crama, Sourour Elloumi, Am{\'e}lie Lambert, and Elisabeth Rodr{\'i}guez-Heck.
\newblock ``Quadratization and convexification in polynomial binary optimization''.
\newblock Journal of Combinatorial Optimization {\bf 50}, Article 28~(2025).

\bibitem{Herrman2021}
Rebekah Herrman, Lorna Treffert, James Ostrowski, Phillip~C. Lotshaw, Travis~S. Humble, and George Siopsis.
\newblock ``Globally optimizing QAOA circuit depth for constrained optimization problems''.
\newblock Algorithms {\bf 14}, 294~(2021).
\newblock arXiv:2108.03281.

\bibitem{Schmidbauer2024}
Lukas Schmidbauer, Karen Wintersperger, Elisabeth Lobe, and Wolfgang Mauerer.
\newblock ``Polynomial reduction methods and their impact on QAOA circuits''.
\newblock In 2024 IEEE International Conference on Quantum Software (QSW), pages 35--45~(2024).
\newblock doi:10.1109/QSW62656.2024.00018.
\newblock arXiv:2406.08889.

\bibitem{Schmidbauer2026}
Lukas Schmidbauer, Elisabeth Lobe, Ina Schaefer, and Wolfgang Mauerer.
\newblock ``It's quick to be square: Fast quadratisation for quantum toolchains''.
\newblock ACM Transactions on Quantum Computing {\bf 7}, Article 18, 1--46~(2026).
\newblock arXiv:2411.19934.

\bibitem{Rovara2026}
Damian Rovara, Lukas Burgholzer, and Robert Wille.
\newblock ``Quantum hardware-efficient selection of auxiliary variables for QUBO formulations''.
\newblock In 2026 Design, Automation \& Test in Europe Conference (DATE), pages 1--7~(2026).
\newblock doi:10.23919/DATE69613.2026.11539407.
\newblock arXiv:2511.19613.

\bibitem{Egger2021}
Daniel~J. Egger, Jakub Mare{\v c}ek, and Stefan Woerner.
\newblock ``Warm-starting quantum optimization''.
\newblock Quantum {\bf 5}, 479~(2021).

\bibitem{TateSDP2023}
Reuben Tate, Majid Farhadi, Creston Herold, Greg Mohler, and Swati Gupta.
\newblock ``Bridging classical and quantum with SDP initialized warm-starts for QAOA''.
\newblock ACM Transactions on Quantum Computing {\bf 4}, Article 9, 1--39~(2023).

\bibitem{TateCustom2023}
Reuben Tate, Jai Moondra, Bryan Gard, Greg Mohler, and Swati Gupta.
\newblock ``Warm-started QAOA with custom mixers provably converges and computationally beats Goemans--Williamson's Max-Cut at low circuit depths''.
\newblock Quantum {\bf 7}, 1121~(2023).

\bibitem{WangMaxCut2018}
Zhihui Wang, Stuart Hadfield, Zhang Jiang, and Eleanor G. Rieffel.
\newblock ``Quantum approximate optimization algorithm for MaxCut: A fermionic view''.
\newblock Physical Review A {\bf 97}, 022304~(2018).


\bibitem{FarhiLocal2020}
Edward Farhi, David Gamarnik, and Sam Gutmann.
\newblock ``The Quantum Approximate Optimization Algorithm Needs to See the Whole Graph: Worst Case Examples''.
\newblock Preprint at arXiv:2005.08747~(2020).


\bibitem{Akshay2020}
V. Akshay, H. Philathong, M. E. S. Morales, and J. D. Biamonte.
\newblock ``Reachability Deficits in Quantum Approximate Optimization''.
\newblock Physical Review Letters {\bf 124}, 090504~(2020).


\bibitem{Cain2023}
Madelyn Cain, Edward Farhi, Sam Gutmann, Daniel Ranard, and Eugene Tang.
\newblock ``The QAOA gets stuck starting from a good classical string''.
\newblock Preprint at arXiv:2207.05089v2~(2023).


\bibitem{LandDoig1960}
A. H. Land and A. G. Doig.
\newblock ``An Automatic Method of Solving Discrete Programming Problems''.
\newblock Econometrica {\bf 28}, 497--520~(1960).


\bibitem{LawlerWood1966}
E. L. Lawler and D. E. Wood.
\newblock ``Branch-and-Bound Methods: A Survey''.
\newblock Operations Research {\bf 14}, 699--719~(1966).


\bibitem{Hadfield2019}
Stuart Hadfield, Zhihui Wang, Bryan O'Gorman, Eleanor G. Rieffel, Davide Venturelli, and Rupak Biswas.
\newblock ``From the Quantum Approximate Optimization Algorithm to a Quantum Alternating Operator Ansatz''.
\newblock Algorithms {\bf 12}, 34~(2019).


\bibitem{Powell1994}
M. J. D. Powell.
\newblock ``A Direct Search Optimization Method That Models the Objective and Constraint Functions by Linear Interpolation''.
\newblock In Advances in Optimization and Numerical Analysis, pages 51--67~(1994).


\bibitem{Student1908}
Student.
\newblock ``The Probable Error of a Mean''.
\newblock Biometrika {\bf 6}, 1--25~(1908).


\bibitem{Lucas2014}
Andrew Lucas.
\newblock ``Ising formulations of many NP problems''.
\newblock Frontiers in Physics {\bf 2}, 5~(2014).


\bibitem{Glover2018}
Fred Glover, Gary Kochenberger, and Yu Du.
\newblock ``A Tutorial on Formulating and Using QUBO Models''.
\newblock Preprint at arXiv:1811.11538~(2018).


\bibitem{CowtanPhase2020}
Alexander Cowtan, Silas Dilkes, Ross Duncan, Will Simmons, and Seyon Sivarajah.
\newblock ``Phase Gadget Synthesis for Shallow Circuits''.
\newblock Electronic Proceedings in Theoretical Computer Science {\bf 318}, 213--228~(2020).


\bibitem{CowtanUCC2020}
Alexander Cowtan, Will Simmons, and Ross Duncan.
\newblock ``A Generic Compilation Strategy for the Unitary Coupled Cluster Ansatz''.
\newblock Preprint at arXiv:2007.10515~(2020).


\bibitem{Murali2019}
Prakash Murali, Jonathan M. Baker, Ali Javadi-Abhari, Frederic T. Chong, and Margaret Martonosi.
\newblock ``Noise-Adaptive Compiler Mappings for Noisy Intermediate-Scale Quantum Computers''.
\newblock Proceedings of the Twenty-Fourth International Conference on Architectural Support for Programming Languages and Operating Systems, 1015--1029~(2019).


\bibitem{DebEtAl2002}
Kalyanmoy Deb, Amrit Pratap, Sameer Agarwal, and T.~Meyarivan.
\newblock ``A fast and elitist multiobjective genetic algorithm: NSGA-II''.
\newblock IEEE Transactions on Evolutionary Computation {\bf 6}, 182--197~(2002).

\bibitem{Basso2022}
Joao Basso, David Gamarnik, Song Mei, and Leo Zhou.
\newblock ``Performance and limitations of the QAOA at constant levels on large sparse hypergraphs and spin glass models''.
\newblock In 2022 IEEE 63rd Annual Symposium on Foundations of Computer Science (FOCS), pages 335--343~(2022).
\newblock doi:10.1109/FOCS54457.2022.00039.
\newblock arXiv:2204.10306.

\bibitem{PelofskeShortDepth2024}
Elijah Pelofske, Andreas B{\"a}rtschi, and Stephan Eidenbenz.
\newblock ``Short-depth QAOA circuits and quantum annealing on higher-order Ising models''.
\newblock npj Quantum Information {\bf 10}, 30~(2024).

\bibitem{Pelofske2024}
Elijah Pelofske, Andreas B{\"a}rtschi, Lukasz Cincio, John Golden, and Stephan Eidenbenz.
\newblock ``Scaling whole-chip QAOA for higher-order Ising spin glass models on heavy-hex graphs''.
\newblock npj Quantum Information {\bf 10}, 109~(2024).
\newblock arXiv:2312.00997.

\bibitem{Qiskit2024}
Ali Javadi-Abhari et al.
\newblock ``Quantum computing with Qiskit''.
\newblock Preprint at arXiv:2405.08810~(2024).


\bibitem{Virtanen2020}
Pauli Virtanen et al.
\newblock ``SciPy 1.0: fundamental algorithms for scientific computing in Python''.
\newblock Nature Methods {\bf 17}, 261--272~(2020).


\bibitem{Li2019}
Gushu Li, Yufei Ding, and Yuan Xie.
\newblock ``Tackling the Qubit Mapping Problem for NISQ-Era Quantum Devices''.
\newblock Proceedings of the Twenty-Fourth International Conference on Architectural Support for Programming Languages and Operating Systems, 1001--1014~(2019).


\bibitem{Crooks2018}
Gavin E. Crooks.
\newblock ``Performance of the Quantum Approximate Optimization Algorithm on the Maximum Cut Problem''.
\newblock Preprint at arXiv:1811.08419~(2018).


\bibitem{Sack2021}
Stefan H. Sack and Maksym Serbyn.
\newblock ``Quantum annealing initialization of the quantum approximate optimization algorithm''.
\newblock Quantum {\bf 5}, 491~(2021).


\bibitem{HuangfuHall2018}
Q. Huangfu and J. A. J. Hall.
\newblock ``Parallelizing the dual revised simplex method''.
\newblock Mathematical Programming Computation {\bf 10}, 119--142~(2018).


\bibitem{Marshall2020}
Jeffrey Marshall, Filip Wudarski, Stuart Hadfield, and Tad Hogg.
\newblock ``Characterizing local noise in QAOA circuits''.
\newblock IOP SciNotes {\bf 1}, 025208~(2020).


\bibitem{WangNoise2021}
Samson Wang et al.
\newblock ``Noise-induced barren plateaus in variational quantum algorithms''.
\newblock Nature Communications {\bf 12}, 6961~(2021).


\bibitem{Franca2021}
Daniel Stilck França and Raul García-Patrón.
\newblock ``Limitations of optimization algorithms on noisy quantum devices''.
\newblock Nature Physics {\bf 17}, 1221--1227~(2021).


\bibitem{Nagies2025}
Sebastian Nagies, Kevin~T. Geier, Javed Akram, Dimitrios Bantounas, Michael Johanning, and Philipp Hauke.
\newblock ``Boosting quantum annealing performance through direct polynomial unconstrained binary optimization''.
\newblock Quantum Science and Technology {\bf 10}, 035008~(2025).
\newblock arXiv:2412.04398.

\bibitem{Nusslein2023}
Jonas N{\"u}{\ss}lein, Sebastian Zielinski, Thomas Gabor, Claudia Linnhoff-Popien, and Sebastian Feld.
\newblock ``Solving (Max) 3-SAT via quadratic unconstrained binary optimization''~(2023).
\newblock arXiv:2302.03536.

\bibitem{Wurtz2021}
Jonathan Wurtz and Peter Love.
\newblock ``MaxCut quantum approximate optimization algorithm performance guarantees for $p>1$''.
\newblock Physical Review A {\bf 103}, 042612~(2021).


\bibitem{BassoHighDepth2022}
Joao Basso, Edward Farhi, Kunal Marwaha, Benjamin Villalonga, and Leo Zhou.
\newblock ``The Quantum Approximate Optimization Algorithm at High Depth for MaxCut on Large-Girth Regular Graphs and the Sherrington-Kirkpatrick Model''.
\newblock Leibniz International Proceedings in Informatics (TQC 2022) {\bf 232}, 7:1--7:21~(2022).


\bibitem{SheraliAdams1990}
Hanif D. Sherali and Warren P. Adams.
\newblock ``A Hierarchy of Relaxations between the Continuous and Convex Hull Representations for Zero-One Programming Problems''.
\newblock SIAM Journal on Discrete Mathematics {\bf 3}, 411--430~(1990).


\bibitem{McCormick1976}
Garth P. McCormick.
\newblock ``Computability of global solutions to factorable nonconvex programs: Part I --- Convex underestimating problems''.
\newblock Mathematical Programming {\bf 10}, 147--175~(1976).


% -------------------- Added in revision --------------------

\bibitem{Koch2026}
Frederik Koch, Shahram Panahiyan, Rick Mukherjee, Joseph Doetsch, et al.
\newblock ``Resource-efficient quantum optimization via higher-order encoding''.
\newblock EPJ Quantum Technology~(2026).
\newblock arXiv:2511.17545.

\bibitem{Mirkarimi2024}
Puya Mirkarimi, David C. Hoyle, Ross Williams, and Nicholas Chancellor.
\newblock ``Quantum optimization with linear Ising penalty functions for customer data science''.
\newblock New Journal of Physics {\bf 26}, 103005~(2024).

\bibitem{Sriluckshmy2023}
P.~V. Sriluckshmy, Vicente Pina-Canelles, Mario Ponce, Manuel~G. Algaba, Fedor {\v S}imkovic, and Martin Leib.
\newblock ``Optimal, hardware native decomposition of parameterized multi-qubit Pauli gates''.
\newblock Quantum Science and Technology {\bf 8}, 045029~(2023).

\bibitem{Amy2018}
Matthew Amy, Parsiad Azimzadeh, and Michele Mosca.
\newblock ``On the CNOT-complexity of CNOT-PHASE circuits''.
\newblock Quantum Science and Technology {\bf 4}, 015002~(2018).

\bibitem{Holmes2020}
Adam Holmes, Sonika Johri, Gian~Giacomo Guerreschi, James~S. Clarke, and A.~Y. Matsuura.
\newblock ``Impact of qubit connectivity on quantum algorithm performance''.
\newblock Quantum Science and Technology {\bf 5}, 025009~(2020).

\bibitem{Zulehner2019}
Alwin Zulehner, Alexandru Paler, and Robert Wille.
\newblock ``An efficient methodology for mapping quantum circuits to the IBM QX architectures''.
\newblock IEEE Transactions on Computer-Aided Design of Integrated Circuits and Systems {\bf 38}, 1226--1236~(2019).

\bibitem{Tan2021}
Benjamin Tan, Marc-Antoine Lemonde, Supanut Thanasilp, Jirawat Tangpanitanon, and Dimitris~G. Angelakis.
\newblock ``Qubit-efficient encoding schemes for binary optimisation problems''.
\newblock Quantum {\bf 5}, 454~(2021).

\bibitem{Herrman2021depth}
Rebekah Herrman, James Ostrowski, Travis~S. Humble, and George Siopsis.
\newblock ``Lower bounds on circuit depth of the quantum approximate optimization algorithm''.
\newblock Quantum Information Processing {\bf 20}, 59~(2021).

\bibitem{Abbas2024}
Amira Abbas et al.
\newblock ``Challenges and opportunities in quantum optimization''.
\newblock Nature Reviews Physics {\bf 6}, 718--735~(2024).

\bibitem{Ayodele2022}
Mayowa Ayodele.
\newblock ``Penalty weights in QUBO formulations: Permutation problems''.
\newblock In Evolutionary Computation in Combinatorial Optimization (EvoCOP 2022), Lecture Notes in Computer Science {\bf 13222}, pages 159--174~(2022).

\end{thebibliography}
)
%
% The entries below were carried over from the original manuscript
% (quantumarticle format). IEEE TQE uses numbered references ordered by
% first citation. Because the list is typeset with thebibliography rather
% than BibTeX, the \bibitem order must match citation order in the final
% text and each entry must be reformatted to IEEE reference style, e.g.
%   E. Boros and P. L. Hammer, "Pseudo-Boolean optimization," Discrete Appl.
%   Math., vol. 123, no. 1-3, pp. 155-225, 2002.
% Newly added references requested in review appear at the end of the list
% under "Added in revision". All new entries must be verified against the
% published versions before submission.
% =====================================================================
\begin{thebibliography}{99}

\bibitem{BorosHammer2002}
Endre Boros and Peter~L. Hammer.
\newblock ``Pseudo-boolean optimization''.
\newblock Discrete Applied Mathematics {\bf 123}, 155--225~(2002).

\bibitem{Hadfield2021}
Stuart Hadfield.
\newblock ``On the representation of Boolean and real functions as Hamiltonians for quantum computing''.
\newblock ACM Transactions on Quantum Computing {\bf 2}, Article 18, 1--21~(2021).
\newblock arXiv:1804.09130.

\bibitem{Farhi2014}
Edward Farhi, Jeffrey Goldstone, and Sam Gutmann.
\newblock ``A quantum approximate optimization algorithm''~(2014).
\newblock arXiv:1411.4028.

\bibitem{Zhou2020}
Leo Zhou, Sheng-Tao Wang, Soonwon Choi, Hannes Pichler, and Mikhail D. Lukin.
\newblock ``Quantum Approximate Optimization Algorithm: Performance, Mechanism, and Implementation on Near-Term Devices''.
\newblock Physical Review X {\bf 10}, 021067~(2020).


\bibitem{Glos2022}
Adam Glos, Aleksandra Krawiec, and Zolt{\'a}n Zimbor{\'a}s.
\newblock ``Space-efficient binary optimization for variational quantum computing''.
\newblock npj Quantum Information {\bf 8}, 39~(2022).
\newblock arXiv:2009.07309.

\bibitem{CampbellDahl2022}
Colin Campbell and Edward Dahl.
\newblock ``QAOA of the highest order''.
\newblock In 2022 IEEE 19th International Conference on Software Architecture Companion (ICSA-C), pages 141--146~(2022).
\newblock doi:10.1109/ICSA-C54293.2022.00035.
\newblock arXiv:2111.12754.

\bibitem{Stein2023}
Jonas Stein, Farbod Chamanian, Maximilian Zorn, Jonas N{\"u}{\ss}lein, Sebastian Zielinski, Michael K{\"o}lle, and Claudia Linnhoff-Popien.
\newblock ``Evidence that PUBO outperforms QUBO when solving continuous optimization problems with the QAOA''.
\newblock In Proceedings of the Genetic and Evolutionary Computation Conference Companion.
\newblock Pages 2254--2262~(2023).
\newblock arXiv:2305.03390.

\bibitem{Bell2025}
Kristina Bell, Adam Lowe, Emily Coles, Nathanael Ridgway, and Gillian Marshall.
\newblock ``A comparison of quadratic and higher-order representations for QAOA''~(2025).
\newblock arXiv:2509.20127.

\bibitem{Dragoi2026}
Sabina Dr{\u a}goi, Alberto Baiardi, and Daniel~J. Egger.
\newblock ``Approximate quadratization of high-order Hamiltonians for combinatorial quantum optimization''.
\newblock Physical Review Research {\bf 8}, 023159~(2026).
\newblock arXiv:2505.04700.

\bibitem{Verchere2023}
Zo{\'e} Verch{\`e}re, Sourour Elloumi, and Andrea Simonetto.
\newblock ``Optimizing variational circuits for higher-order binary optimization''.
\newblock In 2023 IEEE International Conference on Quantum Computing and Engineering (QCE), volume 1, pages 19--25~(2023).
\newblock doi:10.1109/QCE57702.2023.00011.
\newblock arXiv:2307.16756.

\bibitem{Weidenfeller2022}
Johannes Weidenfeller, Lucia~C. Valor, Julien Gacon, Caroline Tornow, Luciano Bello, Stefan Woerner, and Daniel~J. Egger.
\newblock ``Scaling of the quantum approximate optimization algorithm on superconducting qubit based hardware''.
\newblock Quantum {\bf 6}, 870~(2022).
\newblock arXiv:2202.03459.

\bibitem{Harrigan2021}
Matthew P. Harrigan et al.
\newblock ``Quantum approximate optimization of non-planar graph problems on a planar superconducting processor''.
\newblock Nature Physics {\bf 17}, 332--336~(2021).


\bibitem{Rosenberg1975}
I.~G. Rosenberg.
\newblock ``Reduction of bivalent maximization to the quadratic case''.
\newblock Cahiers du Centre d'\'Etudes de Recherche Op\'erationnelle {\bf 17}, 71--74~(1975).

\bibitem{BorosGruber2014}
Endre Boros and Aritanan Gruber.
\newblock ``On quadratization of pseudo-boolean functions''~(2014).
\newblock arXiv:1404.6538.

\bibitem{Dattani2019}
Nike Dattani.
\newblock ``Quadratization in discrete optimization and quantum mechanics''~(2019).
\newblock arXiv:1901.04405.

\bibitem{FreedmanDrineas2005}
Daniel Freedman and Petros Drineas.
\newblock ``Energy minimization via graph cuts: Settling what is possible''.
\newblock In 2005 IEEE Computer Society Conference on Computer Vision and Pattern Recognition (CVPR'05), volume 2, pages 939--946~(2005).
\newblock doi:10.1109/CVPR.2005.143.

\bibitem{Ishikawa2011}
Hiroshi Ishikawa.
\newblock ``Transformation of general binary MRF minimization to the first-order case''.
\newblock IEEE Transactions on Pattern Analysis and Machine Intelligence {\bf 33}, 1234--1249~(2011).

\bibitem{AnthonyEtAl2017}
Martin Anthony, Endre Boros, Yves Crama, and Aritanan Gruber.
\newblock ``Quadratic reformulations of nonlinear binary optimization problems''.
\newblock Mathematical Programming {\bf 162}, 115--144~(2017).

\bibitem{Anthony2016}
Martin Anthony, Endre Boros, Yves Crama, and Aritanan Gruber.
\newblock ``Quadratization of symmetric pseudo-Boolean functions''.
\newblock Discrete Applied Mathematics {\bf 203}, 1--12~(2016).


\bibitem{BorosCramaRodriguezHeck2020}
Endre Boros, Yves Crama, and Elisabeth Rodr{\'i}guez-Heck.
\newblock ``Compact quadratizations for pseudo-Boolean functions''.
\newblock Journal of Combinatorial Optimization {\bf 39}, 687--707~(2020).

\bibitem{VermaLewis2020}
Amit Verma and Mark Lewis.
\newblock ``Optimal quadratic reformulations of fourth degree pseudo-Boolean functions''.
\newblock Optimization Letters {\bf 14}, 1557--1569~(2020).

\bibitem{VermaLewisKochenberger2021}
Amit Verma, Mark Lewis, and Gary Kochenberger.
\newblock ``Efficient QUBO transformation for higher degree pseudo Boolean functions''~(2021).
\newblock arXiv:2107.11695.

\bibitem{CramaEtAl2025}
Yves Crama, Sourour Elloumi, Am{\'e}lie Lambert, and Elisabeth Rodr{\'i}guez-Heck.
\newblock ``Quadratization and convexification in polynomial binary optimization''.
\newblock Journal of Combinatorial Optimization {\bf 50}, Article 28~(2025).

\bibitem{Herrman2021}
Rebekah Herrman, Lorna Treffert, James Ostrowski, Phillip~C. Lotshaw, Travis~S. Humble, and George Siopsis.
\newblock ``Globally optimizing QAOA circuit depth for constrained optimization problems''.
\newblock Algorithms {\bf 14}, 294~(2021).
\newblock arXiv:2108.03281.

\bibitem{Schmidbauer2024}
Lukas Schmidbauer, Karen Wintersperger, Elisabeth Lobe, and Wolfgang Mauerer.
\newblock ``Polynomial reduction methods and their impact on QAOA circuits''.
\newblock In 2024 IEEE International Conference on Quantum Software (QSW), pages 35--45~(2024).
\newblock doi:10.1109/QSW62656.2024.00018.
\newblock arXiv:2406.08889.

\bibitem{Schmidbauer2026}
Lukas Schmidbauer, Elisabeth Lobe, Ina Schaefer, and Wolfgang Mauerer.
\newblock ``It's quick to be square: Fast quadratisation for quantum toolchains''.
\newblock ACM Transactions on Quantum Computing {\bf 7}, Article 18, 1--46~(2026).
\newblock arXiv:2411.19934.

\bibitem{Rovara2026}
Damian Rovara, Lukas Burgholzer, and Robert Wille.
\newblock ``Quantum hardware-efficient selection of auxiliary variables for QUBO formulations''.
\newblock In 2026 Design, Automation \& Test in Europe Conference (DATE), pages 1--7~(2026).
\newblock doi:10.23919/DATE69613.2026.11539407.
\newblock arXiv:2511.19613.

\bibitem{Egger2021}
Daniel~J. Egger, Jakub Mare{\v c}ek, and Stefan Woerner.
\newblock ``Warm-starting quantum optimization''.
\newblock Quantum {\bf 5}, 479~(2021).

\bibitem{TateSDP2023}
Reuben Tate, Majid Farhadi, Creston Herold, Greg Mohler, and Swati Gupta.
\newblock ``Bridging classical and quantum with SDP initialized warm-starts for QAOA''.
\newblock ACM Transactions on Quantum Computing {\bf 4}, Article 9, 1--39~(2023).

\bibitem{TateCustom2023}
Reuben Tate, Jai Moondra, Bryan Gard, Greg Mohler, and Swati Gupta.
\newblock ``Warm-started QAOA with custom mixers provably converges and computationally beats Goemans--Williamson's Max-Cut at low circuit depths''.
\newblock Quantum {\bf 7}, 1121~(2023).

\bibitem{WangMaxCut2018}
Zhihui Wang, Stuart Hadfield, Zhang Jiang, and Eleanor G. Rieffel.
\newblock ``Quantum approximate optimization algorithm for MaxCut: A fermionic view''.
\newblock Physical Review A {\bf 97}, 022304~(2018).


\bibitem{FarhiLocal2020}
Edward Farhi, David Gamarnik, and Sam Gutmann.
\newblock ``The Quantum Approximate Optimization Algorithm Needs to See the Whole Graph: Worst Case Examples''.
\newblock Preprint at arXiv:2005.08747~(2020).


\bibitem{Akshay2020}
V. Akshay, H. Philathong, M. E. S. Morales, and J. D. Biamonte.
\newblock ``Reachability Deficits in Quantum Approximate Optimization''.
\newblock Physical Review Letters {\bf 124}, 090504~(2020).


\bibitem{Cain2023}
Madelyn Cain, Edward Farhi, Sam Gutmann, Daniel Ranard, and Eugene Tang.
\newblock ``The QAOA gets stuck starting from a good classical string''.
\newblock Preprint at arXiv:2207.05089v2~(2023).


\bibitem{LandDoig1960}
A. H. Land and A. G. Doig.
\newblock ``An Automatic Method of Solving Discrete Programming Problems''.
\newblock Econometrica {\bf 28}, 497--520~(1960).


\bibitem{LawlerWood1966}
E. L. Lawler and D. E. Wood.
\newblock ``Branch-and-Bound Methods: A Survey''.
\newblock Operations Research {\bf 14}, 699--719~(1966).


\bibitem{Hadfield2019}
Stuart Hadfield, Zhihui Wang, Bryan O'Gorman, Eleanor G. Rieffel, Davide Venturelli, and Rupak Biswas.
\newblock ``From the Quantum Approximate Optimization Algorithm to a Quantum Alternating Operator Ansatz''.
\newblock Algorithms {\bf 12}, 34~(2019).


\bibitem{Powell1994}
M. J. D. Powell.
\newblock ``A Direct Search Optimization Method That Models the Objective and Constraint Functions by Linear Interpolation''.
\newblock In Advances in Optimization and Numerical Analysis, pages 51--67~(1994).


\bibitem{Student1908}
Student.
\newblock ``The Probable Error of a Mean''.
\newblock Biometrika {\bf 6}, 1--25~(1908).


\bibitem{Lucas2014}
Andrew Lucas.
\newblock ``Ising formulations of many NP problems''.
\newblock Frontiers in Physics {\bf 2}, 5~(2014).


\bibitem{Glover2018}
Fred Glover, Gary Kochenberger, and Yu Du.
\newblock ``A Tutorial on Formulating and Using QUBO Models''.
\newblock Preprint at arXiv:1811.11538~(2018).


\bibitem{CowtanPhase2020}
Alexander Cowtan, Silas Dilkes, Ross Duncan, Will Simmons, and Seyon Sivarajah.
\newblock ``Phase Gadget Synthesis for Shallow Circuits''.
\newblock Electronic Proceedings in Theoretical Computer Science {\bf 318}, 213--228~(2020).


\bibitem{CowtanUCC2020}
Alexander Cowtan, Will Simmons, and Ross Duncan.
\newblock ``A Generic Compilation Strategy for the Unitary Coupled Cluster Ansatz''.
\newblock Preprint at arXiv:2007.10515~(2020).


\bibitem{Murali2019}
Prakash Murali, Jonathan M. Baker, Ali Javadi-Abhari, Frederic T. Chong, and Margaret Martonosi.
\newblock ``Noise-Adaptive Compiler Mappings for Noisy Intermediate-Scale Quantum Computers''.
\newblock Proceedings of the Twenty-Fourth International Conference on Architectural Support for Programming Languages and Operating Systems, 1015--1029~(2019).


\bibitem{DebEtAl2002}
Kalyanmoy Deb, Amrit Pratap, Sameer Agarwal, and T.~Meyarivan.
\newblock ``A fast and elitist multiobjective genetic algorithm: NSGA-II''.
\newblock IEEE Transactions on Evolutionary Computation {\bf 6}, 182--197~(2002).

\bibitem{Basso2022}
Joao Basso, David Gamarnik, Song Mei, and Leo Zhou.
\newblock ``Performance and limitations of the QAOA at constant levels on large sparse hypergraphs and spin glass models''.
\newblock In 2022 IEEE 63rd Annual Symposium on Foundations of Computer Science (FOCS), pages 335--343~(2022).
\newblock doi:10.1109/FOCS54457.2022.00039.
\newblock arXiv:2204.10306.

\bibitem{PelofskeShortDepth2024}
Elijah Pelofske, Andreas B{\"a}rtschi, and Stephan Eidenbenz.
\newblock ``Short-depth QAOA circuits and quantum annealing on higher-order Ising models''.
\newblock npj Quantum Information {\bf 10}, 30~(2024).

\bibitem{Pelofske2024}
Elijah Pelofske, Andreas B{\"a}rtschi, Lukasz Cincio, John Golden, and Stephan Eidenbenz.
\newblock ``Scaling whole-chip QAOA for higher-order Ising spin glass models on heavy-hex graphs''.
\newblock npj Quantum Information {\bf 10}, 109~(2024).
\newblock arXiv:2312.00997.

\bibitem{Qiskit2024}
Ali Javadi-Abhari et al.
\newblock ``Quantum computing with Qiskit''.
\newblock Preprint at arXiv:2405.08810~(2024).


\bibitem{Virtanen2020}
Pauli Virtanen et al.
\newblock ``SciPy 1.0: fundamental algorithms for scientific computing in Python''.
\newblock Nature Methods {\bf 17}, 261--272~(2020).


\bibitem{Li2019}
Gushu Li, Yufei Ding, and Yuan Xie.
\newblock ``Tackling the Qubit Mapping Problem for NISQ-Era Quantum Devices''.
\newblock Proceedings of the Twenty-Fourth International Conference on Architectural Support for Programming Languages and Operating Systems, 1001--1014~(2019).


\bibitem{Crooks2018}
Gavin E. Crooks.
\newblock ``Performance of the Quantum Approximate Optimization Algorithm on the Maximum Cut Problem''.
\newblock Preprint at arXiv:1811.08419~(2018).


\bibitem{Sack2021}
Stefan H. Sack and Maksym Serbyn.
\newblock ``Quantum annealing initialization of the quantum approximate optimization algorithm''.
\newblock Quantum {\bf 5}, 491~(2021).


\bibitem{HuangfuHall2018}
Q. Huangfu and J. A. J. Hall.
\newblock ``Parallelizing the dual revised simplex method''.
\newblock Mathematical Programming Computation {\bf 10}, 119--142~(2018).


\bibitem{Marshall2020}
Jeffrey Marshall, Filip Wudarski, Stuart Hadfield, and Tad Hogg.
\newblock ``Characterizing local noise in QAOA circuits''.
\newblock IOP SciNotes {\bf 1}, 025208~(2020).


\bibitem{WangNoise2021}
Samson Wang et al.
\newblock ``Noise-induced barren plateaus in variational quantum algorithms''.
\newblock Nature Communications {\bf 12}, 6961~(2021).


\bibitem{Franca2021}
Daniel Stilck França and Raul García-Patrón.
\newblock ``Limitations of optimization algorithms on noisy quantum devices''.
\newblock Nature Physics {\bf 17}, 1221--1227~(2021).


\bibitem{Nagies2025}
Sebastian Nagies, Kevin~T. Geier, Javed Akram, Dimitrios Bantounas, Michael Johanning, and Philipp Hauke.
\newblock ``Boosting quantum annealing performance through direct polynomial unconstrained binary optimization''.
\newblock Quantum Science and Technology {\bf 10}, 035008~(2025).
\newblock arXiv:2412.04398.

\bibitem{Nusslein2023}
Jonas N{\"u}{\ss}lein, Sebastian Zielinski, Thomas Gabor, Claudia Linnhoff-Popien, and Sebastian Feld.
\newblock ``Solving (Max) 3-SAT via quadratic unconstrained binary optimization''~(2023).
\newblock arXiv:2302.03536.

\bibitem{Wurtz2021}
Jonathan Wurtz and Peter Love.
\newblock ``MaxCut quantum approximate optimization algorithm performance guarantees for $p>1$''.
\newblock Physical Review A {\bf 103}, 042612~(2021).


\bibitem{BassoHighDepth2022}
Joao Basso, Edward Farhi, Kunal Marwaha, Benjamin Villalonga, and Leo Zhou.
\newblock ``The Quantum Approximate Optimization Algorithm at High Depth for MaxCut on Large-Girth Regular Graphs and the Sherrington-Kirkpatrick Model''.
\newblock Leibniz International Proceedings in Informatics (TQC 2022) {\bf 232}, 7:1--7:21~(2022).


\bibitem{SheraliAdams1990}
Hanif D. Sherali and Warren P. Adams.
\newblock ``A Hierarchy of Relaxations between the Continuous and Convex Hull Representations for Zero-One Programming Problems''.
\newblock SIAM Journal on Discrete Mathematics {\bf 3}, 411--430~(1990).


\bibitem{McCormick1976}
Garth P. McCormick.
\newblock ``Computability of global solutions to factorable nonconvex programs: Part I --- Convex underestimating problems''.
\newblock Mathematical Programming {\bf 10}, 147--175~(1976).


% -------------------- Added in revision --------------------

\bibitem{Koch2026}
Frederik Koch, Shahram Panahiyan, Rick Mukherjee, Joseph Doetsch, et al.
\newblock ``Resource-efficient quantum optimization via higher-order encoding''.
\newblock EPJ Quantum Technology~(2026).
\newblock arXiv:2511.17545.

\bibitem{Mirkarimi2024}
Puya Mirkarimi, David C. Hoyle, Ross Williams, and Nicholas Chancellor.
\newblock ``Quantum optimization with linear Ising penalty functions for customer data science''.
\newblock New Journal of Physics {\bf 26}, 103005~(2024).

\bibitem{Sriluckshmy2023}
P.~V. Sriluckshmy, Vicente Pina-Canelles, Mario Ponce, Manuel~G. Algaba, Fedor {\v S}imkovic, and Martin Leib.
\newblock ``Optimal, hardware native decomposition of parameterized multi-qubit Pauli gates''.
\newblock Quantum Science and Technology {\bf 8}, 045029~(2023).

\bibitem{Amy2018}
Matthew Amy, Parsiad Azimzadeh, and Michele Mosca.
\newblock ``On the CNOT-complexity of CNOT-PHASE circuits''.
\newblock Quantum Science and Technology {\bf 4}, 015002~(2018).

\bibitem{Holmes2020}
Adam Holmes, Sonika Johri, Gian~Giacomo Guerreschi, James~S. Clarke, and A.~Y. Matsuura.
\newblock ``Impact of qubit connectivity on quantum algorithm performance''.
\newblock Quantum Science and Technology {\bf 5}, 025009~(2020).

\bibitem{Zulehner2019}
Alwin Zulehner, Alexandru Paler, and Robert Wille.
\newblock ``An efficient methodology for mapping quantum circuits to the IBM QX architectures''.
\newblock IEEE Transactions on Computer-Aided Design of Integrated Circuits and Systems {\bf 38}, 1226--1236~(2019).

\bibitem{Tan2021}
Benjamin Tan, Marc-Antoine Lemonde, Supanut Thanasilp, Jirawat Tangpanitanon, and Dimitris~G. Angelakis.
\newblock ``Qubit-efficient encoding schemes for binary optimisation problems''.
\newblock Quantum {\bf 5}, 454~(2021).

\bibitem{Herrman2021depth}
Rebekah Herrman, James Ostrowski, Travis~S. Humble, and George Siopsis.
\newblock ``Lower bounds on circuit depth of the quantum approximate optimization algorithm''.
\newblock Quantum Information Processing {\bf 20}, 59~(2021).

\bibitem{Abbas2024}
Amira Abbas et al.
\newblock ``Challenges and opportunities in quantum optimization''.
\newblock Nature Reviews Physics {\bf 6}, 718--735~(2024).

\bibitem{Ayodele2022}
Mayowa Ayodele.
\newblock ``Penalty weights in QUBO formulations: Permutation problems''.
\newblock In Evolutionary Computation in Combinatorial Optimization (EvoCOP 2022), Lecture Notes in Computer Science {\bf 13222}, pages 159--174~(2022).

\end{thebibliography}


% =====================================================================
% biography.tex  --  author biographies for main.tex (input by % =====================================================================
% biography.tex  --  author biographies for main.tex (input by % =====================================================================
% biography.tex  --  author biographies for main.tex (input by \input{biography})
%
% TQE biography conventions (from the official template):
%   * Photograph 1 in wide by 1.25 in high, at least 300 dpi, placed in photos/.
%   * First paragraph: place and date of birth (optional), degrees with field,
%     institution, city, country, and year.
%   * Second paragraph: employment history, current position, research interests.
%   * Third paragraph: memberships, awards, and editorial roles.
%   * Use IEEEbiographynophoto if a photograph is unavailable.
% Order must match the author list on the first page.
% Replace each placeholder with the text supplied by the author.
% =====================================================================

\begin{IEEEbiography}[{\includegraphics[width=1in,height=1.25in,clip,keepaspectratio]{placeholder.png}}]{Jingcong Liu}
is with the Department of Mathematics, University of Warwick, Coventry, U.K.
Biography to be supplied.
\end{IEEEbiography}

\begin{IEEEbiography}[{\includegraphics[width=1in,height=1.25in,clip,keepaspectratio]{placeholder.png}}]{Yiding Su}
is with the Department of Mathematics, University of Warwick, Coventry, U.K.
Biography to be supplied.
\end{IEEEbiography}

\begin{IEEEbiography}[{\includegraphics[width=1in,height=1.25in,clip,keepaspectratio]{placeholder.png}}]{Ruojun Zhang}
is with the School of Electronics and Computer Science, University of Southampton,
Southampton, U.K. Biography to be supplied.
\end{IEEEbiography}

\begin{IEEEbiography}[{\includegraphics[width=1in,height=1.25in,clip,keepaspectratio]{placeholder.png}}]{Hafiz Muhammad Waseem}
is with the Warwick Manufacturing Group, University of Warwick, Coventry, U.K.
Biography to be supplied.
\end{IEEEbiography}

\begin{IEEEbiography}[{\includegraphics[width=1in,height=1.25in,clip,keepaspectratio]{placeholder.png}}]{Bo Chen}
is with the Warwick Business School, University of Warwick, Coventry, U.K.
Biography to be supplied.
\end{IEEEbiography}

\begin{IEEEbiography}[{\includegraphics[width=1in,height=1.25in,clip,keepaspectratio]{placeholder.png}}]{Carsten Maple}
is with the Warwick Manufacturing Group, University of Warwick, Coventry, U.K.
Biography to be supplied.
\end{IEEEbiography}
)
%
% TQE biography conventions (from the official template):
%   * Photograph 1 in wide by 1.25 in high, at least 300 dpi, placed in photos/.
%   * First paragraph: place and date of birth (optional), degrees with field,
%     institution, city, country, and year.
%   * Second paragraph: employment history, current position, research interests.
%   * Third paragraph: memberships, awards, and editorial roles.
%   * Use IEEEbiographynophoto if a photograph is unavailable.
% Order must match the author list on the first page.
% Replace each placeholder with the text supplied by the author.
% =====================================================================

\begin{IEEEbiography}[{\includegraphics[width=1in,height=1.25in,clip,keepaspectratio]{placeholder.png}}]{Jingcong Liu}
is with the Department of Mathematics, University of Warwick, Coventry, U.K.
Biography to be supplied.
\end{IEEEbiography}

\begin{IEEEbiography}[{\includegraphics[width=1in,height=1.25in,clip,keepaspectratio]{placeholder.png}}]{Yiding Su}
is with the Department of Mathematics, University of Warwick, Coventry, U.K.
Biography to be supplied.
\end{IEEEbiography}

\begin{IEEEbiography}[{\includegraphics[width=1in,height=1.25in,clip,keepaspectratio]{placeholder.png}}]{Ruojun Zhang}
is with the School of Electronics and Computer Science, University of Southampton,
Southampton, U.K. Biography to be supplied.
\end{IEEEbiography}

\begin{IEEEbiography}[{\includegraphics[width=1in,height=1.25in,clip,keepaspectratio]{placeholder.png}}]{Hafiz Muhammad Waseem}
is with the Warwick Manufacturing Group, University of Warwick, Coventry, U.K.
Biography to be supplied.
\end{IEEEbiography}

\begin{IEEEbiography}[{\includegraphics[width=1in,height=1.25in,clip,keepaspectratio]{placeholder.png}}]{Bo Chen}
is with the Warwick Business School, University of Warwick, Coventry, U.K.
Biography to be supplied.
\end{IEEEbiography}

\begin{IEEEbiography}[{\includegraphics[width=1in,height=1.25in,clip,keepaspectratio]{placeholder.png}}]{Carsten Maple}
is with the Warwick Manufacturing Group, University of Warwick, Coventry, U.K.
Biography to be supplied.
\end{IEEEbiography}
)
%
% TQE biography conventions (from the official template):
%   * Photograph 1 in wide by 1.25 in high, at least 300 dpi, placed in photos/.
%   * First paragraph: place and date of birth (optional), degrees with field,
%     institution, city, country, and year.
%   * Second paragraph: employment history, current position, research interests.
%   * Third paragraph: memberships, awards, and editorial roles.
%   * Use IEEEbiographynophoto if a photograph is unavailable.
% Order must match the author list on the first page.
% Replace each placeholder with the text supplied by the author.
% =====================================================================

\begin{IEEEbiography}[{\includegraphics[width=1in,height=1.25in,clip,keepaspectratio]{placeholder.png}}]{Jingcong Liu}
is with the Department of Mathematics, University of Warwick, Coventry, U.K.
Biography to be supplied.
\end{IEEEbiography}

\begin{IEEEbiography}[{\includegraphics[width=1in,height=1.25in,clip,keepaspectratio]{placeholder.png}}]{Yiding Su}
is with the Department of Mathematics, University of Warwick, Coventry, U.K.
Biography to be supplied.
\end{IEEEbiography}

\begin{IEEEbiography}[{\includegraphics[width=1in,height=1.25in,clip,keepaspectratio]{placeholder.png}}]{Ruojun Zhang}
is with the School of Electronics and Computer Science, University of Southampton,
Southampton, U.K. Biography to be supplied.
\end{IEEEbiography}

\begin{IEEEbiography}[{\includegraphics[width=1in,height=1.25in,clip,keepaspectratio]{placeholder.png}}]{Hafiz Muhammad Waseem}
is with the Warwick Manufacturing Group, University of Warwick, Coventry, U.K.
Biography to be supplied.
\end{IEEEbiography}

\begin{IEEEbiography}[{\includegraphics[width=1in,height=1.25in,clip,keepaspectratio]{placeholder.png}}]{Bo Chen}
is with the Warwick Business School, University of Warwick, Coventry, U.K.
Biography to be supplied.
\end{IEEEbiography}

\begin{IEEEbiography}[{\includegraphics[width=1in,height=1.25in,clip,keepaspectratio]{placeholder.png}}]{Carsten Maple}
is with the Warwick Manufacturing Group, University of Warwick, Coventry, U.K.
Biography to be supplied.
\end{IEEEbiography}


\EOD

\end{document}
